\selectlanguage{ngerman}%
\frenchspacing%
\begin{ledgroupsized}[r]{120mm}%
\footnotesize%
\pstart%
\noindent%
\textbf{Überlieferung:}%
\pend%
\end{ledgroupsized}%
\begin{ledgroupsized}[r]{114mm}%
\footnotesize%
\pstart%
\parindent -6mm%
\makebox[6mm][l]{\textit{L}}%
Notiz:
LH XXXVII\;4 Bl.\;95.
Ein Zettel (10\;x\;6,5 cm);
Ränder beschnitten;
Papier\-erhal\-tungs\-maßnahmen.
Eine Seite auf Bl.\;\href{\Scan{LH_37_04_095r}}{95\,r\textsuperscript{o}}.
Bl.\;\href{\Scan{LH_37_04_095v}}{95\,v\textsuperscript{o}} ist unbeschrieben.
\pend%
\end{ledgroupsized}%
\vspace{5mm}%
\begin{ledgroup}%
\footnotesize%
\pstart%
\noindent%
\textbf{Datierungsgründe:}
Der Träger der undatierten Notiz N.\;\ref{60250}, die eine Bemerkung zum Ursprung der Körperkohäsion festhält, weist weder Spuren eines Wasserzeichens noch sonstige datierungsrelevante äußere Merkmale auf. %
Der Text zeigt aber eine bemerkenswerte terminologische Nähe zur \textit{Hypothesis physica nova} (Mainz 1671).\cite{00256} %
So etwa ist der Ausdruck \textit{aether, qui est fortasse ille spiritus Domini} (unten, \refpassage{LH_37_04_095_095r_spiritus-1}{LH_37_04_095_095r_spiritus-2}) nahezu identisch in der wohl im Winter 1670/1671 entstandenen Abhandlung anzutreffen:
\textit{AEHTER (is enim fortasse est ille spiritus Dominis qui super aquis ferebatur ...} (§\;7; \textit{LSB} VI,\;2 N.\;40, S.\;225.9\,f.).\cite{00256}
Und auch die Darstelluung der durch die Bewegung des Äthers "animierten" Blasen (\textit{bullae}) \textit{tanquam rerum stamina} (unten, \refpassage{LH_37_04_095_095r_stamina-1}{LH_37_04_095_095r_stamina-2}) findet ihr Pendant in der frühen naturphilosophischen Schrift:
\textit{Hae jam bullae sunt semina rerum, stamina specierum, receptacula aetheris, corporum basis, consistentiae causa et fundamentum tantae varietatis quantam in rebus, tanti impetus quantum in motibus admiramus} (§\;12; \textit{LSB} VI,\;2 N.\;40, S.\;226.14--16).\cite{00256}
Die Annahme liegt somit nahe, dass die Notiz N.\;\ref{60250} in einem ähnlichen Zeitraum wie die \textit{Hypothesis physica nova}\cite{00256} entstand, d.\,h. um den Winter 1670/1671.
\pend%
\pstart%
Es gilt jedoch zu beachten, dass Leibniz in seinen physikalischen Überlegungen über die Ursache der Körperkohäsion noch recht lange auf die "Blasen"-Begrifflichkeit zurückgriff, so etwa noch in dem auf die Zeitspanne zwischen Ende Januar und März\,/\,April 1683 datierbaren Entwurf \textit{Explicatio mechanica elastri} (\textit{LSB} VIII,\;3 N.\;14\textsubscript{3}, S.\;206--208).\cite{01437}
Und der Gedanke, dass von der Bewegung eines dünnen Fluidums wie des Äthers der Zusammenhalt der Körper herrühre, lässt sich sogar bis in die frühen Jahre des 18.\;Jahrhundert verfolgen (siehe etwa \textit{LSB} VIII,\;3 N.\;33).\cite{02608}
Daher ist nicht auszuschließen, dass die Notiz N.\;\ref{60250} bis wenigstens zum Frühjahr 1683 entstanden sein könnte, als die "Blasen"-Begrifflichkeit noch in Verwendung war.
Die terminologische Nähe zur \textit{Hypothesis physica nova}\cite{00256} lässt aber eine frühere Entstehung vermuten.%
\pend%
\end{ledgroup}%
\selectlanguage{latin}%
\frenchspacing%
\vspace{8mm}%
\normalsize%
\pstart%
\noindent%
\lbrack%
\href{\Scan{LH_37_04_095r}}%
{95\,r\textsuperscript{o}}%
\rbrack\
Quod
\edtext{%
connexio corporum%
\indexS{connexio corporum}%
\indexS{corpora connexa}%
}{%
\lemma{connexio}%
\Bfootnote{%
\textit{(1)}~rerum\indexS{connexio rerum} %
\textit{(2)}~corporum%
~\textit{L}%
}}
oriatur a motu%
\indexS{motus origo connexionis}%
\indexS{origo connexionis corporum}
confirmat alabastri pulvis,%
\indexS{pulvis alabastri}%
\indexS{alabaster}
qui fusus ebullit,%
\indexS{pulvis ebullens}
bullae autem connexionem partium habent.%
\indexS{bulla habens connexionem partium}%
\indexS{connexio partium}%
\indexS{partes connexae}
\edtext{%
Cessante calore%
\indexS{calor cessans}%
}{%
\lemma{Cessante}%
\Bfootnote{%
\textit{(1)}~ebullitione\indexS{ebullitio} %
\textit{(2)}~calore%
~\textit{L}%
}}
reperitur merus pulvis.
Eodem modo
\edlabel{LH_37_04_095_095r_spiritus-1}%
aether,%
\indexS{aether interlabens}
qui est fortasse ille spiritus Domini%
\indexS{spiritus Domini}%
\edlabel{LH_37_04_095_095r_spiritus-2}%
\lbrack,\rbrack\
motu suo%
\indexS{motus aetheris}
massam%
\indexS{massa}
quae
ipsius interlabentis actione sublata%
\indexS{actio aetheris}%
\indexS{aether interlabens}
in pulverem mortuum concideret,%
\indexS{pulvis mortuus}
\edlabel{LH_37_04_095_095r_stamina-1}%
in bullas %
tanquam rerum stamina%
\indexS{bullae tanquam stamina rerum}%
\indexS{stamina rerum}
animavit.%
\edlabel{LH_37_04_095_095r_stamina-2}
Idem confirmat limatura ferri,%
\indexS{limatura ferri}%
\indexS{ferrum}
quae
interlabente magnetis admoti effluvio%
\indexS{effluvius magnetis}%
\indexS{magnes}
quasi in filamenta concrescit.%
\indexS{filamentum limaturae}
\pend%
\selectlanguage{latin}%
\frenchspacing%
\normalsize%
\newpage%
